\scanpage{348}
\mspage{\unsure{124|$12_4$}}

\deleted{\unsure{cheminer par des}}

\deleted{\unsure{cheminer|cheminons} par des considérations \inserted{$\pm$} générales.

Je suppose que l’on ait dans $U$ une sous-catégorie $V$ ouverte \deleted{ou fermée (elle sera fermée dans le cas qui nous intéresse)}, isomorphe à $\unsure{\Upsilon|Y}$ ; plus précisément $V$ est \deleted{\unsure{relative à}} \inserted{associée à} un point $t_0\in\dot U$ \deleted{comme} par
\[
\begin{aligned}
\text{\textcircled{a}}\quad\biggl\{
&\deleted{\unsure{V=U_{/t}}}\ \unsure{V}=\unsure{U}_{/t_0}\ \deleted{\unsure{et}}\quad\text{si }t\in\dot U\setminus\{t_0\}\\
&\deleted{\unsure{V_{/t_0}\longrightarrow V}\quad\unsure{V_{\!/t}}\ \text{\unsure{est d’un type}}}\ \unsure{\text{(ordinal }1\text{)}|\text{(élément }1\text{)}}.
\end{aligned}
\]
\unsure{(} le foncteur canonique $U_{/t_0}\unsure{\text{---}|\longrightarrow}U$ étant une immersion ouverte (ce qui est automatique si $\overline U$, voire $X$ tout entier, est une catégorie ordonnée).

Je m’aperçois à l’instant que dans le cas qui m’intéresse particulièrement, \deleted{cette} \inserted{la} catégorie \inserted{$U_{/t}$ elle-$=$,} n’est pas \inserted{\unsure{isom à}} $\Upsilon$, \inserted{comme je croyais l’avoir \unsure{lu|vu},} mais la catégorie un peu plus grande
\[
\deleted{\unsure{\Phi_2}}
\]
\diagram{\begin{tikzcd}[ampersand replacement=\&,column sep=0.9em,row sep=tiny]
\Phi_2=\bigl\{ \& {}\arrow[r] \& {} \& {}\arrow[l]\arrow[r] \& {} \& {}\arrow[l] \& \bigr\},\\
{} \& {} \& {} \& \underbrace{\hspace{3.5em}}_{\Upsilon}
\end{tikzcd}}
dans laquelle $\Upsilon$ est immergée en position centrale stratégique. \deleted{Néanmoins, je vais supposer \unsure{que $V$} \illegible{3}} Le point, \deleted{\unsure{crucial}} est d’utiliser seulement \inserted{la structure intrinsèque de} $V$, c’est \inserted{\textcircled{b} $V$ est asphérique et} \textcircled{b} qu’il existe une \deleted{\unsure{fonction}} \inserted{flèche}
\[
\deleted{\unsure{V}}\qquad p_V^0:V\longrightarrow V'_0
\]
dans Cat, telle que
\[
p_{V!}^0(e_V)=e_{V'_0},\qquad V'_0\ \uline{\text{non asphérique}}.
\]
Soit donc $V'$ déduite de $V'_0$ en lui ajoutant un élément final, \unsure{d’où|donc} $V'_0\hookrightarrow\unsure{V'_0|V'}$ \unsure{(imm.|l’imm.} ouverte),}

% archive sheet number: 348
